\par\bigskip\pdfbookmark[0]{On the Theory of Biquadratic Residues. IV.}{nachlass-biquadratic-IV}\centerline{\textbf{[IV.]}}\medskip
\renewcommand{\Article}[1]{\par\medskip\Needspace{4\baselineskip}\pdfbookmark[1]{Art. #1}{nachlass-biquadratic-IV-art-#1}\centerline{\textbf{#1.}}\nobreak\smallskip}
\Article{1}
\[
\begin{array}{lll}
\text{Biquadratic residue?}&m=a+bi;&aa+bb=d\\[3pt]
\text{Modulus}&M=A+Bi,&AA+BB=D
\end{array}
\]
\[
\frac{mD}{M}=\mu\qquad
\mu=\alpha+\beta i,\quad\alpha=aA+bB,\quad\beta=Ab-Ba
\]
\[
\xi+\eta i=\pi;\quad\pi m=x+yi=p;\quad\pi M=X+Yi=P
\]
Relations
\[
\begin{array}{lll}
x=a\xi-b\eta&d\xi=ax+by&D\xi=AX+BY\\
y=b\xi+a\eta&d\eta=-bx+ay&D\eta=-BX+AY\\
X=A\xi-B\eta&dX=\alpha x+\beta y&Dx=\alpha X-\beta Y\\
Y=B\xi+A\eta&dY=-\beta x+\alpha y&Dy=\beta X+\alpha Y
\end{array}
\]

\SourceLeaf{353}
\[
\begin{aligned}
\beta\xi&=-Bx+bX=Ay-aY\\
\beta\eta&=-Ax+aX=-By+bY\\
\alpha\xi&=Ax+aY=By+aX\\
\alpha\eta&=-Bx+aY=Ay-bX
\end{aligned}
\]

Those $\pi$ for which $\xi$ and $\eta$ lie between $0$ and $\frac12$ shall be denoted by $\pi^0$, and the corresponding $p$ and $P$ by $p^0$ and $P^0$; those $\pi$ for which $\eta=0$ and $\xi$ lies between $0$ and $\frac12$, by $\pi'$; those for which $\xi=\frac12$ and $\eta$ lies between $0$ and $\frac12$, by $\pi''$; those $\pi$ for which $\eta=\frac12$ and $\xi$ lies between $0$ and $\frac12$, by $\pi'''$; finally, those for which $\xi=0$ and $\eta$ lies between $0$ and $\frac12$, by $\pi''''$.

The decident of $\frac{m}{M}$ is found as follows:

Collect all \emph{integer} $P^0$; the corresponding $x^0$ and $y^0$ will therefore be fractional. Let the respective intensors of $p^0$ be $t^0$, that is, the numbers $0,1.\ 2,3$, according as
\[
\begin{array}{lllll}
[x^0]&\text{even,}&\text{odd,}&\text{odd,}&\text{even}\\
\relax[y^0]&\text{even,}&\text{even,}&\text{odd,}&\text{odd}
\end{array}
\]
The required decident is then $=\Sigma\pm t^0$, where the upper sign is to be taken for even $P^0$, the lower for odd.

This is the \emph{first} method.

\Needspace{20\baselineskip}\Article{2}
We now wish to arrange the individual $P^0$ according to the values of $Y^0$. Restricting ourselves to the case in which $a,b,A,B$ are positive, the smallest value is given by $Y^0\dotstwo+1$, the largest by $\frac12(A+B-1)$. For each particular value of $Y^0$, the values of $X^0$ must lie between certain bounds, namely

I. when $A-B$ is positive
\[
\begin{array}{c|c|c}
\text{if}&\text{between}&\text{and}\\
Y<\frac12B&-\frac{BY^0}{A}&\frac{AY^0}{B}\\[5pt]
Y=\frac12B&-\frac{BB}{2A}&\frac12A\\[5pt]
Y>\frac12B\text{ and }<\frac12A&-\frac{BY^0}{A}&-\frac{BY^0}{A}+\frac{D}{2A}\\[5pt]
Y>\frac12A&\frac{AY^0}{B}-\frac{D}{2B^0}&-\frac{BY^0}{A}+\frac{D}{2A}
\end{array}
\]

\SourceLeaf{354}
II. When $A-B$ is negative,
\[
\begin{array}{c|c|c}
\text{if}&\text{between}&\text{and}\\
Y<\frac12A&-\frac{BY^0}{A}&\frac{AY^0}{B}\\[5pt]
Y>\frac12A\text{ and }<\frac12B&\frac{AY^0}{B}-\frac{D}{2B}&\frac{AY^0}{B}\\[5pt]
Y=\frac12B&\frac{AB-D}{2B}&\frac12A\\[5pt]
Y>\frac12B&\frac{AY^0}{B}-\frac{D}{2B}&-\frac{BY^0}{A}+\frac{D}{2A}
\end{array}
\]

At the lower bounds either $\xi=0$ or $\eta=\frac12$; at the upper bounds, on the other hand, either $\eta=0$ or $\xi=\frac12$. It is easy to prove that the bounds of $x$ are never integers.

\Needspace{20\baselineskip}\Article{3}
We now wish to express the partial sums for each particular $Y^0$ in another way. At the bounds $p$ will have particular values, which may be denoted by $p^*,p^{**}$, and while $X$ changes continuously from one bound to the other, $p$ will pass continuously from $p^*$ to $p^{**}$. But the integer contained in $[p]$ changes here in jumps, with either the real or the imaginary part always changing by one unit. The changes occur at the values of $X$
\[
X',\ X'',\ X'''\ldots X^\mu
\]
which are already ordered by size and to which the values of $p$
\[
p',\ p''.\ p'''\ldots
\]
correspond.

The last $X^\mu$ can also be identical with $X^{**}$ when $\beta$ is positive, or $X'$ identical with $X^*$, and so on.

The $x$ are increasing here; thus, when $x^{**}$ is an integer, it is counted as $x^\mu$.

The $y$ are increasing for positive $\beta$; then $y^{**}$ is fully included in the count;
\begin{center}decreasing for negative $\beta$; then $y^*$ is included in the count.\end{center}

\SourceLeaf{355}
Let the intensors of $p^*,p^{**}$ be $\lambda^*$ and $\lambda^{**}$
\[
\begin{array}{rll}
\text{the intensor from }p'&\text{up to }p''\dotstwo&\lambda^*+\delta'\\
p''&\text{up to }p'''\dotstwo&\lambda^*+\delta'+\delta''\\
p^\mu&\text{up to }p^{**}\dotstwo&\lambda^*+\delta'+\delta''\dotstwo+\delta^\mu=\lambda^{**}
\end{array}
\]
Then the partial sum, insofar as $Y^0$ is even, becomes
\[
\begin{aligned}
={}&\lambda^*(g-h)+(\lambda^*+\delta')(g'-h')\\
&+(\lambda^*+\delta'+\delta'')(g''-h'')+\text{etc.}+\lambda^{**}(g^\mu-h^\mu)
\end{aligned}
\]
where $g$ denotes the number of even $X^0$ from $X^*$ to $X'$, and $h$ that of the odd ones.

\clearpage\Article{4}
This formula can also be expressed as follows:
\[
\begin{aligned}
&\lambda^*\{[\tfrac12X^*-\tfrac12]-[\tfrac12X^*]\}\\
&+\delta'\{[\tfrac12X'-\tfrac12]-[\tfrac12X']\}\\
&+\delta''\{[\tfrac12X''-\tfrac12]-[\tfrac12X'']\}\\
&+\text{etc.}\\
&+\delta^\mu\{[\tfrac12X^\mu-\tfrac12]-[\tfrac12X^\mu]\}\\
&-\lambda^{**}\{[\tfrac12X^{**}-\tfrac12]-[\tfrac12X^{**}]\}
\end{aligned}
\]
or by
\[
\lambda^*\varepsilon^*+\delta'\varepsilon'+\delta''\varepsilon''+\text{etc.}+\delta^\mu\varepsilon^\mu-\lambda^{**}\varepsilon^{**}
\]
\[
\begin{array}{rll}
\text{where in general }\varepsilon=0&\text{if }[X]\text{ odd}\\
\text{and }=-1&\text{if }[X]\text{ is even and }Y\text{ even}\\
+1&\text{if }[X]\text{ even and }Y\text{ odd.}
\end{array}
\]

For $\delta$, on the other hand, the values are
\[
\begin{array}{ll|rr}
&&\multicolumn{2}{c}{\beta}\\
&&\text{positive}&\text{negative}\\
\text{if }x\text{ an even integer,}&[y]\text{ even}&-1&-1\\
&[y]\text{ odd}&+1&+1\\
x\text{ an odd integer,}&[y]\text{ even}&+1&+1\\
&[y]\text{ odd}&-1&-1\\
y\text{ an even integer,}&[x]\text{ even}&-3&+3\\
&[x]\text{ odd}&-1&+1\\
y\text{ an odd integer,}&[x]\text{ even}&+3&-3\\
&[x]\text{ odd}&+1&-1
\end{array}
\]

\SourceLeaf{356}
\clearpage\Article{5}
From this we derive the following second method of determining the decident.

I. Collect all combinations of integer values of $Y$ and $x$ having the following properties:
\begin{enumerate}
\item that $\xi=\frac{Ax+bY}{\alpha}$ lies between $0$ and $\frac12$, with the second bound included;
\item that $\eta=\frac{-Bx+aY}{\alpha}$ lies between $0$ and $\frac12$, with the first bound included.
\end{enumerate}

II. For these calculate
\[
\begin{aligned}
X&=\frac{Dx+\beta Y}{\alpha}\\[3pt]
y&=\frac{\beta x+dY}{\alpha}
\end{aligned}
\]

III. Omit all those for which $[X]$ is odd, and divide the remaining ones, for which $[X]$ is even, into two classes;

in the first class put those for which, simultaneously,
\[
\begin{array}{l}
Y\text{ even, }x\text{ even, }[y]\text{ even}\\
\text{or where one of these conditions holds;}
\end{array}
\]
in the second class, those for which two of these conditions, or none at all, hold,
\[
\begin{array}{rl}
\text{or in I.}&\text{where }[Y+x+y]\text{ even}\\
\text{II.}&\text{where }[Y+x+y]\text{ odd}
\end{array}
\]
and call the excess of the number in the first class over that in the second $c$.

IV. Collect all combinations of integer values of $Y$ and $y$ having the following properties:
\begin{enumerate}
\item that $\xi=\frac{Ay-aY}{\beta}$ lies between $0$ and $\frac12$, with the first bound included when $\beta$ is negative, the second included when $\beta$ is positive;
\item that $\eta=\frac{-By+bY}{\beta}$ lies between $0$ and $\frac12$, with the first bound included for positive $\beta$, the second for negative.
\end{enumerate}
\BookRule{.3\textwidth}
